\documentclass{article}
\usepackage[left=2cm, right=2cm, top=2cm, bottom=2cm]{geometry}
\usepackage{graphicx}
\usepackage{amsmath}
\usepackage{amssymb}
\usepackage{amsthm}
\usepackage{fancyhdr}
\usepackage{verbatim}
\usepackage{listings}
\usepackage{xcolor}
\usepackage{pdfpages}
\usepackage{cancel}
\usepackage{physics}
\usepackage{esint}
\usepackage{braket}

\lstset{
    language=C++,
    basicstyle=\ttfamily\footnotesize,
    keywordstyle=\color{blue},
    commentstyle=\color{green},
    stringstyle=\color{red},
    numbers=left,
    numberstyle=\tiny\color{gray},
    frame=single,
    breaklines=true,
    captionpos=b,
}

\title{PHYS 496: Lecture \# 1}
\author{Cliff Sun}

\newtheorem{theorem}{Theorem}[section]
\newtheorem{lemma}[theorem]{Lemma}
\newtheorem{definition}[theorem]{Definition}
\newtheorem{conjecture}[theorem]{Conjecture}
\newtheorem{proposition}[theorem]{Proposition}
\newtheorem{corollary}[theorem]{Corollary}
\newtheorem{one minute paper}[theorem]{One Minute Paper}

\pagestyle{fancy}
\lhead{\textbf{\thepage}\ \ \nouppercase{\rightmark}}
\chead{PHYS 496: Lecture \# 1}
\rhead{Cliff Sun}

\begin{document}

\maketitle

\section*{Lecture Span}
\begin{enumerate}
    \item Introduction to class
    \item How to read physics papers
    \item Writing effective titles
\end{enumerate}

\section*{Introduction to class}

\begin{enumerate}
    \item Homework is due on gradescope
    \item Email phys496@physics.illinois.edu if missing a class
    \item Two colloquium reports due. One due near september. 2nd due in early november. 
\end{enumerate}

\section*{How to read physics papers}

The 4 i's of paper reading. 

\begin{enumerate}
    \item Importance: does the paper contain relevant information for your research? Read title and abstract. Also, look at author list. The goal of every step should be to figure out if you're going to abandon this paper.
    \item Iterate: Find main points of each section, also generate questions while reading. REQUIRES ACTIVE BRAIN HOURS.
    Also, figuring out how to understand a paper is good, like figuring out how the result of the paper changes if you change the underlying model, etc.    
    \item Interpret: write down key ideas in your own words.
    \item Integrate: Figuring out how the information in the paper fits with your current knowledge.  
\end{enumerate}

Explaining something simple is the last step. Writing down things imprints it into the brain. 

\section*{Writing effective titles}

\begin{enumerate}
    \item Make it interesting, but not too interesting. Keep them short. 
    \item Don't use qualitative words
    \item Don't use names of poople, places, etc.
    \item Front load key words
\end{enumerate}

\end{document}